
\medskip
\phantomsection\label{lem:sga3-VIB-7.5}%
\noindent\textbf{Lemma 7.5.---}
\emph{Let \(S\) be a scheme, let \(G\) be an \(S\)-group
scheme, let \(X\) be an \(S\)-prescheme, and let
\(f:X\to G\) be an \(S\)-morphism. Suppose that \(X\) and \(G\)
are locally of finite presentation over \(S\), and that, for
every \(s\in S\) and every maximal point \(g\) of \(G_s\),
there exists a point \(x\) of \(X\) such that \(f(x)=g\) and
\(f\) is flat at \(x\). Then the morphism}
\phantomsection\label{form:sga3-VIB-7.5-covering-morphism}%
\[
  \mu\circ(f\times_S f):
  X\times_S X\longrightarrow G
\]
\emph{is a covering for the \((\mathrm{fppf})\) topology.}

\footnote{\label{note:sga3-VIB-75}Editor's note: the argument
in the original has been expanded in what follows.}
By EGA~IV\(_3\), 11.3.1, the set \(V\) of points of \(X\) at
which \(f\) is flat is open, and \(f|_V\) is an open morphism.
Thus
\[
  U=f(V)
\]
is an open subset of \(G\). Moreover, by hypothesis,
\[
  U\cap G_s
\]
is dense in \(G_s\) for every \(s\in S\). Let \(\phi\) denote
the restriction of \(\mu\circ(f\times_S f)\) to
\(V\times_S V\).

It is enough to show that \(\phi\) is a covering for the
\((\mathrm{fppf})\) topology, and for this it is enough to show
that \(\phi\) is faithfully flat and of finite presentation
(\hyperref[sga3:iv:proposition:6:3:1:item:iv:01]
{Expos\'e~IV, 6.3.1\textup{(iv)}}). Now \(\phi\) is the
composite
\phantomsection\label{form:sga3-VIB-7.5-factorization}%
\[
  V\times_S V
  \xrightarrow{\ f|_V\times f|_V\ }
  U\times_S U
  \xrightarrow{\ \mu\ }
  G,
\]
where the first morphism is faithfully flat and locally of
finite presentation, since \(f|_V\) is. It therefore remains
only to prove the same assertions for the restriction of
\(\mu\) to \(U\times_S U\).

Since \(G\to S\) is locally of finite presentation and flat, so
is
\[
  \mu:G\times_S G\longrightarrow G.
\]
Indeed, this morphism is isomorphic to the morphism obtained
from \(G\to S\) by the base change \(G\to S\)
(\hyperref[sga3:VIA:0.1]{VI\(_A\), 0.1}). The induced morphism
\[
  U\times_S U\longrightarrow G
\]
therefore has the same properties. To prove that this last
morphism is surjective, it is enough to check fiber by fiber.
There the result follows from
\hyperref[sga3:VIA:0.5]{VI\(_A\), 0.5}, since
\(U\cap G_s\) is a dense open subset of \(G_s\) for every
\(s\in S\).
